\documentclass[10pt,letterpaper]{article}
\usepackage{cornell-notes}

\title{Clause 07.05.08.04: Compound Requirements}
\author{}
\date{}
\setDocTitle{Clause 07.05.08.04: Compound Requirements}
\setDocSubtitle{C++ 2024}
\setDocOwner{C++ 2024 Cornell Notes}
\setDocDate{\today}
\setCornellCollection{C++ 2024 Cornell Notes}
\setCornellUnitType{Clause}
\setCornellUnitNumber{07.05.08.04}
\setCornellUnitTitle{Compound Requirements}
\hypersetup{pdftitle={Clause 07.05.08.04: Compound Requirements},pdfsubject={Cornell notes},pdfauthor={}}

\begin{document}
\maketitle
\begin{CornellOverviewBox}{Overview}
\textbf{Classification:} Core language - Normative\\[2pt]
\textbf{Learning objective:} Understand the rules, terminology, and semantic consequences associated with Compound requirements.
\end{CornellOverviewBox}\begin{CornellSummaryBox}{Summary}
{\sffamily\bfseries\color{Accent} Summary}\par\vspace{3pt}Section 7.5.8.4 focuses on Compound requirements. Use the listed grammar productions together with the semantic constraints and disambiguation rules referenced by the section. When applying it, preserve the distinction between mandatory requirements, recommendations, permissions, and behavior deliberately left undefined, unspecified, or implementation-defined.
\end{CornellSummaryBox}

\end{document}
